\documentclass[10pt,a4paper]{amsart}
\usepackage[margin=19mm]{geometry}
\usepackage{amsmath,amssymb,mathtools}
\usepackage[scr=rsfs,cal=euler]{mathalfa}
\usepackage{xfrac}
\usepackage[unicode,hidelinks,bookmarks=false]{hyperref}
\newcommand{\ie}{i.e.\ }
\newcommand{\cf}{cf.\ }
\newcommand{\resp}{respectively\ }
\newcommand{\Subsection}{\subsection}
\providecommand{\nobreakdash}{\nobreak\hbox{-}}
\setlength{\emergencystretch}{2em}
\multlinegap0pt
\usepackage{bm}
\usepackage{bbm}
\usepackage{dsfont}
\usepackage{xspace}
\usepackage{cancel}
\usepackage{mathtools}
\usepackage{amsthm}
\makeatletter
\@ifundefined{thm}{\newtheorem{thm}{Theorem}}{}
\makeatother
\title[Fixed Points of Meir-Keeler and Leader Contractions with bou]{Fixed Points of Meir-Keeler and Leader Contractions with bounded orbits in b-Metric Spaces}
\author{Hassan Khandani}
\date{Source: arXiv:2506.09074v1 (CC BY 4.0)}
\begin{document}
\maketitle

\noindent\emph{Source and licence.} Fixed Points of Meir-Keeler and Leader Contractions with bounded orbits in b-Metric Spaces. Version arXiv:2506.09074v1 (CC BY 4.0), distributed under the Creative Commons Attribution 4.0 International licence, as recorded in the arXiv licence metadata for this exact version and shown on the abstract page https://arxiv.org/abs/2506.09074v1. Full rights notice: licenses/arxiv-2506.09074-CC-BY-4.0.txt.\par\smallskip

\noindent\textit{Editorial scope.} Counted unit: the Thm whose source label is \texttt{thm\_leader} in the article (source0.tex lines 215--217, source label \texttt{thm\_leader}), delivered together with the article's own complete original proof of it (source0.tex lines 218--322, one \texttt{proof} environment of 105 lines). No mathematics is written, repaired or re-derived: statement and proof are the article's own text line for line. Required context retained uncounted: none. Every definition, display and label of those lines is retained unchanged. Omitted from this package and not redistributed: 1--214, 323--550. Nothing inside a retained excerpt is omitted or altered: every retained body is delivered exactly as the article's lines are written, so no omission receipt of any kind accompanies this package. Citations: the retained excerpts carry no citation command at all, so no bibliography entry of the article is transcribed and no reference is redistributed. Reference adaptation: none was needed and none was made; every reference inside the delivered unit resolves inside it, so no unresolved reference or citation is produced. Printed slips kept literal: no printed word, symbol, hypothesis or number of the counted statement or of its proof was altered. Layout and class: the article's own class is \texttt{birkjour} with packages including enumitem, tikz-cd, fontenc, inputenc, caption, amsmath, amssymb; the wrapper is the lane's pinned \texttt{amsart} editorial wrapper at 10pt on A4, which restates the theorem-like environments the retained text needs and adds the article's own macro definitions that the retained text needs (none), with extra wrapper packages (bm, bbm, dsfont, xspace, cancel, mathtools). The delivered numbering is the unit's own. Assets: no retained span contains a figure environment, a graphics-inclusion command, a PDF-page inclusion, a table or a TikZ picture; no image file is copied into this package, the package declares no asset dependency and no image file is redistributed with it. Licence: version arXiv:2506.09074v1 of this article is distributed under the Creative Commons Attribution 4.0 International licence, as recorded in the arXiv licence metadata for this exact version and shown on the abstract page https://arxiv.org/abs/2506.09074v1. A full rights notice is delivered at licenses/arxiv-2506.09074-CC-BY-4.0.txt. Characters: every retained excerpt is pure ASCII, so the delivered engine, XeLaTeX with its default Latin Modern text font, reports no missing character.
\begin{thm}\label{thm_leader}
Let $(X,\Delta)$ be a complete $b$-metric space with coefficient $s\geq 1$, and let $T:X\to X$ a non-expansive Leader contraction with bounded orbits. Then $T$ has a fixed point. 
\end{thm}

\begin{proof}Let $x\in X$, and $\mathcal{O}(x) =\{T^{n}x\}_{n=0}^{\infty}$ be the orbit of $T$ at $x$. We define $\Sigma$ as ordered paires of $(m,n)$, where $m,n$ are sequences of non-negative integers, as follows: 
\begin{equation}
 \begin{aligned}
    \Sigma = \Big\{ 
        \big( m, n \big) \;\big|\;
           m=\{a_k\}_{k=0}^{\infty}, n=\{b_k\}_{k=0}^{\infty}, &  m(k)\leq n(k)\quad\text{for all }k\in \mathbb{Z}^{+}\\
 & ,\lim_{k\to\infty} m(k) = \lim_{k\to\infty} n(k) = \infty
          \Big\}
  \end{aligned}
\end{equation}
for $(m,n)\in \Sigma, p\in \mathbb{Z}^{+}$, define
\begin{align}\label{sig_inf}
\sigma(m,n,p) = \inf_{k\geq 0}d(T^{m(k)+p}x,T^{n(k)+p}x),
\end{align}
\begin{equation}\label{sig_def}
\sigma(p) = \inf_{(m,n)\in \Sigma}\sigma(m,n,p).
\end{equation}
\begin{equation}\label{theta_sup}
\theta(p) = \sup_{(m,n)\in \Sigma}\sigma(m,n,p).
\end{equation}
Since $T$ has bounded orbits, for all $(m,n)\in \sigma, p\in \mathbb{Z}^{+} $ we have:
\[
0\le \sigma(m,n,p), \sigma(p),\theta(p)<\infty.
\]
 On the other hand, since $T$ is non-expansive, we have
\[
\sigma(p+1)\leq \sigma(p)\quad \text{for all }p\in \mathbb{Z}^{+}.
\]
This shows that there exists $\sigma\geq 0$ such that 
\[
\sigma(p)\to \sigma \text{ as }p\to \infty\quad\text{ ,and }\sigma(p)\geq  \sigma\quad\text{ for all }p\in \mathbb{Z}^{+}.
\]
We claim that $\sigma =0$. If $\sigma>0$, let $\delta>0$ and $r\in \mathbb{Z}^{+}$ such that for all $x,y\in X$:
\begin{equation}\label{leader_delta}
d(x,y)<\sigma+\delta\rightarrow d(T^{r}x,T^{r}y) <  \sigma
\end{equation}
There exists $p_0\in \mathbb{Z}^{+}$ such that
\[
\sigma(p)<\sigma+\delta\quad\text{ for all }p\ge p_0
\]
Let $p\ge p_{0}$, invoning \eqref{sig_def}, there exists $(m,n)\in \Sigma$ such that
\[
\sigma(m,n,p)<\sigma+\delta
\]
by \eqref{sig_inf} we have:
\[
\inf_{k\geq 0}d(T^{m(k)+p}x,T^{n(k)+p}x)<\sigma+\delta.
\]
This implies that there exists $k_0\ge 0$ such that:
\[
d(T^{m(k_0)+p}x,T^{n(k_0)+p}x)<\sigma+\delta.
\]
By \eqref{leader_delta} we have:
\begin{equation}\label{lead_1}
d(T^{m(k_0)+p+r}x,T^{n(k_0)+p+r}x) < \sigma.
\end{equation}
Therefore
\[
\inf_{k\geq 0} d(T^{m(k)+p+r}x,T^{n(k)+p+r}x)< \sigma.
\]
This implies that:
\begin{equation}\label{lead_3}
\sigma(m+r,n+r,p) < \sigma.
\end{equation}
Since $(m+r,n+r)\in \Sigma$,  invoking \eqref{sig_inf} we deduce that:
\begin{equation}
\sigma(p)\leq \sigma(m+r,n+r,p) < \sigma.
\end{equation}
This is a contradiction that  proves our claim that is $\sigma = 0$.\\
Since $T$ is non-expansive for all $(m,n)\in \Sigma$ we have:
\[
\sigma(m,n,p+1)\leq \sigma(m,n,p)
\]
This implies that 
\[
\sup_{(m,n)\in\Sigma}\sigma(m,n,p+1)\leq  \inf_{(m,n)\in\Sigma} \sigma(m,n,p),
\]
therefore:
\[
0\leq \theta(p+1) \leq \sigma(p).
\]
This implies that 
\begin{equation}
\lim_{p\to\infty}\theta(p) = 0
\end{equation}
We claim that the sequence $\{T^{n}x\}$ is Cauchy. If not, there exists $\epsilon>0$ such that $m(k),n(k)\to \infty$ as $k\to\infty$, $m(k)\leq n(k)$ for all $k\geq 0$ such that:
\[
\epsilon\leq d(T^{m(k)}x,T^{n(k)}x)
\]
Let $p\in \mathbb{N}$, without loss of generality we assume that $m(k),n(k)\geq p$ for all $k\geq 0$. We have:
\begin{equation}
\epsilon\leq d(T^{m(k)}x,T^{n(k)}x) =  d(T^{m(k)-p+p}x,T^{n(k)-p+p}x).
\end{equation}
We observe that $(m_1,n_1) = (m-p, n-p)\in \Sigma$, therefore we get:
\begin{align}
\epsilon&\leq \inf_{k\geq 0}d(T^{m(k)-p+p}x,T^{n(k)-p+p}x)\nonumber\\
& = \inf_{k\geq 0}d(T^{m_1(k)+p}x,T^{n_1(k)+p}x)\nonumber\\
&= \sigma(m_1,n_1,p)\leq \theta(p)\nonumber.
\end{align}
Hence, we get
\[
\epsilon\leq \theta(p).
\]
Since $p$ is arbitrary, taking the limit as $p\to \infty$ we get $\epsilon = 0$. This contradiction proves our claim, that is $\{T^{n}x\}$ is a Cauchy sequence. Since $(X,\Delta)$ is a complete metric space it converges to a point $z\in X$. Since $T$ is continuous $Tz=z$ and our proof is complete.
\end{proof}

\end{document}
